Choroplethová mapa výdajů na \ac{APZ} jako podílu \ac{HDP} v~EU27 vizuálně potvrzuje geografickou strukturu „sociální investiční státy vs.~liberální periferie“: severozápadní Evropa (DK, SE, BE, NL, DE) investuje výrazně nad mediánem, zatímco středo- a~jihovýchodní Evropa (CZ, SK, HU, RO) je konzistentně pod mediánem.
Tato geografická distribuce není výsledkem ekonomické nutnosti --- ČR by si vyšší \ac{APZ} výdaje mohla dovolit, \ac{HDP}/obyvatele je vyšší než v~řadě zemí se středně vysokou \ac{APZ} --- ale výsledkem politické volby a~institucionální inerce.
Nízká \ac{APZ} výdaje v~ČR jsou přímou výzvou ke kolektivnímu vyjednávání: pokud stát neposkytuje dostatečný rekvalifikační a~podpůrný rámec, pracovníci jsou více závislí na stabilitě zaměstnaneckého vztahu, a~\ac{KS} se stávají ještě důležitějším nástrojem ochrany než v~zemích s~plně funkční \ac{APZ}.
Zvýšení \ac{APZ} výdajů na \SI{1}{\percent} \ac{HDP} --- v~kombinaci s~erga omnes kolektivním vyjednáváním --- je proto synergickým balíčkem, kde \ac{APZ} posiluje flexisecurity (bezpečnost přechodu) a~\ac{KS} posiluje mzdový standard pro ty, kteří v~zaměstnání zůstávají (srov.~obr.~\ref{fig:lmp_expenditure}).

Choroplethová mapa výdajů na \ac{APZ} jako podílu \ac{HDP} v~EU27 vizuálně potvrzuje geografickou strukturu „sociální investiční státy vs.~liberální periferie“: severozápadní Evropa (DK, SE, BE, NL, DE) investuje výrazně nad mediánem, zatímco středo- a~jihovýchodní Evropa (CZ, SK, HU, RO) je konzistentně pod mediánem.
Tato geografická distribuce není výsledkem ekonomické nutnosti --- ČR by si vyšší \ac{APZ} výdaje mohla dovolit, \ac{HDP}/obyvatele je vyšší než v~řadě zemí se středně vysokou \ac{APZ} --- ale výsledkem politické volby a~institucionální inerce.
Nízká \ac{APZ} výdaje v~ČR jsou přímou výzvou ke kolektivnímu vyjednávání: pokud stát neposkytuje dostatečný rekvalifikační a~podpůrný rámec, pracovníci jsou více závislí na stabilitě zaměstnaneckého vztahu, a~\ac{KS} se stávají ještě důležitějším nástrojem ochrany než v~zemích s~plně funkční \ac{APZ}.
Zvýšení \ac{APZ} výdajů na \SI{1}{\percent} \ac{HDP} --- v~kombinaci s~erga omnes kolektivním vyjednáváním --- je proto synergickým balíčkem, kde \ac{APZ} posiluje flexisecurity (bezpečnost přechodu) a~\ac{KS} posiluje mzdový standard pro ty, kteří v~zaměstnání zůstávají (srov.~obr.~\ref{fig:lmp_expenditure}).

Choroplethová mapa výdajů na \ac{APZ} jako podílu \ac{HDP} v~EU27 vizuálně potvrzuje geografickou strukturu „sociální investiční státy vs.~liberální periferie“: severozápadní Evropa (DK, SE, BE, NL, DE) investuje výrazně nad mediánem, zatímco středo- a~jihovýchodní Evropa (CZ, SK, HU, RO) je konzistentně pod mediánem.
Tato geografická distribuce není výsledkem ekonomické nutnosti --- ČR by si vyšší \ac{APZ} výdaje mohla dovolit, \ac{HDP}/obyvatele je vyšší než v~řadě zemí se středně vysokou \ac{APZ} --- ale výsledkem politické volby a~institucionální inerce.
Nízká \ac{APZ} výdaje v~ČR jsou přímou výzvou ke kolektivnímu vyjednávání: pokud stát neposkytuje dostatečný rekvalifikační a~podpůrný rámec, pracovníci jsou více závislí na stabilitě zaměstnaneckého vztahu, a~\ac{KS} se stávají ještě důležitějším nástrojem ochrany než v~zemích s~plně funkční \ac{APZ}.
Zvýšení \ac{APZ} výdajů na \SI{1}{\percent} \ac{HDP} --- v~kombinaci s~erga omnes kolektivním vyjednáváním --- je proto synergickým balíčkem, kde \ac{APZ} posiluje flexisecurity (bezpečnost přechodu) a~\ac{KS} posiluje mzdový standard pro ty, kteří v~zaměstnání zůstávají (srov.~obr.~\ref{fig:lmp_expenditure}).

Choroplethová mapa výdajů na \ac{APZ} jako podílu \ac{HDP} v~EU27 vizuálně potvrzuje geografickou strukturu „sociální investiční státy vs.~liberální periferie“: severozápadní Evropa (DK, SE, BE, NL, DE) investuje výrazně nad mediánem, zatímco středo- a~jihovýchodní Evropa (CZ, SK, HU, RO) je konzistentně pod mediánem.
Tato geografická distribuce není výsledkem ekonomické nutnosti --- ČR by si vyšší \ac{APZ} výdaje mohla dovolit, \ac{HDP}/obyvatele je vyšší než v~řadě zemí se středně vysokou \ac{APZ} --- ale výsledkem politické volby a~institucionální inerce.
Nízká \ac{APZ} výdaje v~ČR jsou přímou výzvou ke kolektivnímu vyjednávání: pokud stát neposkytuje dostatečný rekvalifikační a~podpůrný rámec, pracovníci jsou více závislí na stabilitě zaměstnaneckého vztahu, a~\ac{KS} se stávají ještě důležitějším nástrojem ochrany než v~zemích s~plně funkční \ac{APZ}.
Zvýšení \ac{APZ} výdajů na \SI{1}{\percent} \ac{HDP} --- v~kombinaci s~erga omnes kolektivním vyjednáváním --- je proto synergickým balíčkem, kde \ac{APZ} posiluje flexisecurity (bezpečnost přechodu) a~\ac{KS} posiluje mzdový standard pro ty, kteří v~zaměstnání zůstávají (srov.~obr.~\ref{fig:lmp_expenditure}).

\input{texparts/python/lmp_expenditure_map}



